\chapter{双群与旋量特征标的重建方法}
\label{app:double-groups}

\chapref{chap:rotation-group}已经解释，半整数自旋态在 $2\pi$ 转动下获得负号，因此普通点群不足以作为自旋量子态的线性表示群。本附录把重点放在“怎样从 $SU(2)$ 表示和普通点群自己重建双群特征标”，并列出最常用的分支规则。这样可以把双群表理解为连续群表示限制到有限群后的结果，而不是另一套独立数字表。

\section{双群作为二重覆盖}

双群表比普通点群表多出来的最重要符号是 $\bar E$。它不是另一个空间恒等操作，而是“在旋量空间中绕任意轴转 $2\pi$”对应的中心群元。对普通三维向量，$E$ 与 $\bar E$ 投影到同一个空间操作；对半整数自旋，二者的表示矩阵却相差负号。因此读双群特征标表时，第一件事就是看某一行在 $\bar E$ 上的特征标：若等于在 $E$ 上的特征标，该表示是单值的；若整体反号，则是自旋型双值表示。这个检查可以在还没理解整张表之前先把普通轨道表示与真正的旋量表示区分开。


对纯转动点群 $P\subset SO(3)$，定义其双群为覆盖映射
\[
p:SU(2)\longrightarrow SO(3)
\]
下的原像
\begin{equation}
P^D=p^{-1}(P).
\label{eq:double-point-preimage}
\end{equation}
若 $|P|=N$，则 $|P^D|=2N$。其中存在中心元
\[
\bar E=-I_2,
\qquad
\bar E^2=E,
\]
它在三维空间中对应 $2\pi$ 转动，因此投影到普通点群后与单位元 $E$ 相同，但在旋量表示中不能与 $E$ 识别。

普通点群的每个元素 $R$ 有两个提升 $u$ 与 $\bar E u=-u$。双群不可约表示分成两类：若
\[
D(\bar E)=+I,
\]
则表示下降为普通点群表示，称为单值表示；若
\[
D(\bar E)=-I,
\]
则为真正的旋量表示或双值表示。\chapref{chap:rotation-group}中整数 $j$ 与半整数 $j$ 的差别正对应这两个情况。

含反演的点群还要给状态附加宇称。因为空间反演与纯转动交换，在常见中心对称双群中可以把旋量表示再乘上一维 $g/u$ 宇称表示，所以同一个旋量不可约表示会出现 $+$/$-$ 或 $g/u$ 两个版本。

\section{\texorpdfstring{$SU(2)$ 特征标公式}{SU(2) 特征标公式}}

自旋 $j$ 不可约表示中，绕任意轴转角 $\theta$ 的本征值为
\[
\ee^{-\ii m\theta},
\qquad m=-j,-j+1,\dots,j.
\]
因此特征标只依赖转角而不依赖转轴：
\begin{align}
\chi_j(\theta)
&=\sum_{m=-j}^{j}\ee^{-\ii m\theta}\notag\\
&=\frac{\sin\bigl((2j+1)\theta/2\bigr)}{\sin(\theta/2)}
=\frac{\sin\bigl((j+\tfrac12)\theta\bigr)}{\sin(\theta/2)}.
\label{eq:su2-character-angle}
\end{align}
在 $\theta\to0$ 极限，$\chi_j(0)=2j+1$；而 $2\pi$ 转动给出
\begin{equation}
\chi_j(2\pi)=(-1)^{2j}(2j+1).
\label{eq:central-double-character}
\end{equation}
所以只要知道双群每个共轭类所对应的空间转角，就已经能写出任何 $SU(2)$ 角动量表示限制到该双群后的特征标。

若双群不可约特征标记为 $\chi_\alpha$，那么 $j$ 表示限制后的重数仍由有限群正交公式给出：
\begin{equation}
m_\alpha^{(j)}
=\frac1{|P^D|}\sum_{g\in P^D}
\overline{\chi_\alpha(g)}\,\chi_j(g).
\label{eq:double-group-multiplicity}
\end{equation}
这就是从“角动量 $j$”转成“晶体双群能带标签”的一般算法。

\section{三个典型双群}

\subsection{\texorpdfstring{$D_2^D$ 与四元数群}{D2 双群与四元数群}}

普通 $D_2$ 由绕 $x,y,z$ 的三个 $\pi$ 转动组成。它们在 $SU(2)$ 中的提升为
\[
U_x=-\ii\sigma_x,
\qquad
U_y=-\ii\sigma_y,
\qquad
U_z=-\ii\sigma_z.
\]
利用 Pauli 代数，
\[
U_x^2=U_y^2=U_z^2=-I_2=\bar E,
\]
且
\[
U_xU_y=U_z,
\qquad
U_yU_x=-U_z=\bar E U_z.
\]
因此 $D_2^D$ 与四元数群 $Q_8$ 同构。$Q_8$ 有四个一维普通表示和一个二维旋量表示，维数平方和为
\[
4\times1^2+2^2=8.
\]
这正是\chapref{chap:rotation-group}讨论“原来四个一维 $D_2$ 轨道表示与自旋 $1/2$ 结合后都进入二维旋量表示”的最简代数模型。

\subsection{轴向群与半整数角动量标签}

对只含一条主轴的群，旋量在 $C_n$ 提升下的本征值最自然写成
\[
\exp\left(-\ii m\frac{2\pi}{n}\right),
\qquad m\in\ZZ+\frac12.
\]
这里 $m$ 与 $m+n$ 在普通空间转动上给出相同相位结构，但在双群中还必须追踪 $2\pi$ 转动的负号。实际晶体表常把 $m$ 与 $-m$ 成对组织成二维旋量表示。若再有镜面或横向二重轴，这两个分量会被相互交换，从而形成不可约二维块。

\subsection{\texorpdfstring{立方双群 $O^D$ 与 $O_h^D$}{立方双群 OD 与 OhD}}

在立方对称中，最常见的旋量不可约表示按凝聚态习惯记为
\[
\Gamma_6\ (2\text{维}),
\qquad
\Gamma_7\ (2\text{维}),
\qquad
\Gamma_8\ (4\text{维}).
\]
把连续转动群的 $j$ 表示限制到 $O^D$，最前几项为
\begin{align*}
j=\frac12&:\quad \Gamma_6,\\
j=\frac32&:\quad \Gamma_8,\\
j=\frac52&:\quad \Gamma_7\oplus\Gamma_8,\\
j=\frac72&:\quad \Gamma_6\oplus\Gamma_7\oplus\Gamma_8.
\end{align*}
维数分别检验为 $2,4,6,8$。加入反演以后，每个表示再带 $g/u$ 或 $+/-$ 宇称，例如 $\Gamma_8^+$、$\Gamma_7^-$。

这些分支不是经验表。以 $j=5/2$ 为例，先由 \eqnrefs{eq:su2-character-angle} 在 $O^D$ 各类上计算 $\chi_{5/2}$，再用 \eqnrefs{eq:double-group-multiplicity} 与 $\Gamma_6,\Gamma_7,\Gamma_8$ 逐行取内积，就得到
\[
D^{(5/2)}\downarrow O^D=\Gamma_7\oplus\Gamma_8.
\]
实际能带程序输出的双群标签，本质上就是这一限制与约化过程。

\section{轨道表示与自旋表示的组合}

若不考虑自旋--轨道耦合，一个轨道态属于普通点群不可约表示 $\Gamma_{\rm orb}$，自旋 $1/2$ 属于二维旋量表示 $\Gamma_{1/2}$。考虑自旋--轨道耦合以后，应在双群中分解
\begin{equation}
\Gamma_{\rm orb}\otimes\Gamma_{1/2}
=\bigoplus_\alpha m_\alpha\Gamma_\alpha^D.
\label{eq:orbital-spin-double-product}
\end{equation}
重数仍由特征标乘法与正交公式得到。

以立方 $O_h$ 为例，$d$ 轨道先在普通点群下分成
\[
E_g\oplus T_{2g}.
\]
各自再与自旋 $1/2$ 相乘，在常用双群记号下有
\[
E_g\otimes\Gamma_6^+\cong\Gamma_8^+,
\]
\[
T_{2g}\otimes\Gamma_6^+\cong\Gamma_7^+\oplus\Gamma_8^+.
\]
这说明“考虑自旋以后能带标签改变”并不是在原来每条能带旁边简单加一个 $\uparrow/\downarrow$；真正的本征态属于双群不可约表示，轨道和自旋通常已经不能分离。

\section{时间反演与双群的关系}

双群解决的是空间转动在旋量上的线性表示问题；时间反演是反酉对称，二者不能混为同一种群操作。对半整数自旋且满足 $\Theta^2=-1$ 的系统，Kramers 定理在时间反演不变动量点强制二重简并。若同时存在反演，对每个一般 $\vb k$，$PT$ 联合作用可把 $\vb k$ 送回自身并保持 Kramers 型二重简并。这个简并的来源是反酉对称，而不是某个二维双群不可约表示本身。实际读能带时必须区分“表示维数要求的简并”和“时间反演要求的简并”。
